\chapter*{Lampiran}
\addcontentsline{toc}{chapter}{Lampiran}

% ============================================================
% Lampiran A: Rubrik Penilaian
% ============================================================
\section*{Lampiran A: Rubrik Penilaian Tugas Praktik Kriptografi}

\begin{table}[htbp]
\centering
\small
\begin{tabular}{|>{\raggedright\arraybackslash}p{3cm}|>{\raggedright\arraybackslash}p{3.5cm}|>{\raggedright\arraybackslash}p{3.5cm}|>{\raggedright\arraybackslash}p{3.5cm}|}
\hline
\textbf{Kriteria} & \textbf{Sangat Baik} & \textbf{Baik} & \textbf{Perlu Perbaikan} \\
\hline
Kebenaran Algoritma & Kode benar, output sesuai teori & Kode benar, minor bug & Kode sebagian benar \\
\hline
Implementasi C & Struktur kode rapi, modular & Struktur cukup baik & Struktur berantakan \\
\hline
Dokumentasi & Lengkap, jelas, komentar bermakna & Cukup lengkap & Kurang lengkap \\
\hline
Pengujian & Test cases lengkap & Beberapa test cases & Minimal testing \\
\hline
Analisis & Analisis mendalam terhadap kelemahan & Analisis memadai & Analisis sederhana \\
\hline
\end{tabular}
\caption{Rubrik Penilaian Tugas Praktik Kriptografi}
\end{table}

% ============================================================
% Lampiran B: Contoh Template Laporan
% ============================================================
\section*{Lampiran B: Contoh Template Laporan Tugas}
\begin{enumerate}
  \item Judul dan Identitas Mahasiswa
  \item Deskripsi Algoritma yang Diimplementasikan
  \item Alur Enkripsi/Dekripsi (diagram atau pseudocode)
  \item Implementasi (cuplikan kode C penting)
  \item Hasil Pengujian dengan Berbagai Input
  \item Analisis Kelemahan dan Keamanan
  \item Refleksi dan Kesimpulan
\end{enumerate}

% ============================================================
% Lampiran C: Glosarium
% ============================================================
\section*{Lampiran C: Glosarium Istilah Kriptografi}
\begin{itemize}
  \item \textbf{Plaintext}: Pesan asli sebelum enkripsi
  \item \textbf{Ciphertext}: Pesan terenkripsi
  \item \textbf{Cipher}: Algoritma untuk enkripsi/dekripsi
  \item \textbf{Key}: Kunci rahasia untuk enkripsi/dekripsi
  \item \textbf{Enkripsi}: Proses mengubah plaintext menjadi ciphertext
  \item \textbf{Dekripsi}: Proses mengubah ciphertext kembali ke plaintext
  \item \textbf{Kriptografi simetris}: Menggunakan kunci yang sama untuk enkripsi dan dekripsi
  \item \textbf{Kriptografi asimetris}: Menggunakan pasangan kunci publik dan privat
  \item \textbf{Fungsi hash}: Fungsi satu arah yang memetakan input ke output dengan panjang tetap
  \item \textbf{Kriptanalisis}: Studi untuk memecahkan cipher tanpa mengetahui kunci
  \item \textbf{Confidentiality}: Kerahasiaan informasi
  \item \textbf{Integrity}: Keutuhan data
  \item \textbf{Authenticity}: Keaslian sumber
\end{itemize}

% ============================================================
% Lampiran D: Referensi Tambahan
% ============================================================
\section*{Lampiran D: Referensi Tambahan}
\begin{itemize}
  \item NIST Cryptographic Standards: \url{https://csrc.nist.gov/}
  \item OpenSSL Documentation: \url{https://www.openssl.org/docs/}
  \item Cryptography Stack Exchange: \url{https://crypto.stackexchange.com/}
  \item The Joy of Cryptography (Mike Rosulek): \url{https://joyofcryptography.com/}
\end{itemize}
